\begin{abstract}
Mundelein Center, on Loyola University Chicago's Lakeshore Campus, holds 54 general-purpose classrooms on twelve
floors but is served by only two elevator banks: four cars that reach floor~9 and two cars that reach floor~14.
Each class changeover produces long lines. We build a discrete-time simulation of one 15-minute passing period
to compare two low-cost behavioral measures, holding backpacks in front and walking to and from floors 2 and~3,
and to test whether running floors 10--14 on a time offset spreads the load. In the baseline, mean waiting time
is 5.2 minutes and 23\% of students arrive late. Backpacks cut mean wait by 31\%; walking cuts it by only 6\%;
together they cut it by 35\% and lateness to 11\%. Offsetting the upper floors shortens the time to clear the
queue by about 20\% but, combined with the other measures, \emph{raises} the share of late students.

These results are conditional on the model in two ways. The parameters are estimates, not measurements. More
consequentially, the simulator embeds fixed behavioral and operational assumptions (for example, that students
never abandon a queue and that the controller treats all cars alike), which the sensitivity analyses cannot test.
We list them, state which conclusions are most exposed, and treat absolute minutes as outputs of one model rather
than facts about Mundelein Center.
\end{abstract}
